% ============================================================
\chapter*{Preface}
\addcontentsline{toc}{chapter}{Preface}
\markboth{Preface}{Preface}
% ============================================================

\section*{Objectives of this course}

Bayesian statistics is one of the two major paradigms of statistical
inference, alongside the frequentist approach. Founded on Bayes'
theorem, it provides a coherent framework for quantifying uncertainty,
incorporating prior knowledge, and updating beliefs in light of data.

This course is aimed at Master-level students in applied mathematics,
data science, and artificial intelligence. It assumes prior mastery of
probability theory, elementary mathematical statistics, and linear
algebra.

\section*{Course organization}

The course is structured into eleven chapters, grouped into four
thematic parts:

\begin{enumerate}
  \item \textbf{Foundations} (Chapters 1--5): Bayesian paradigm,
        conjugate families, estimation, testing, and prediction.
  \item \textbf{Advanced models} (Chapter 6): hierarchical models.
  \item \textbf{Computational methods} (Chapters 7--9): MCMC
        (Metropolis--Hastings, Gibbs sampler) and variational inference.
  \item \textbf{Extensions and applications} (Chapters 10--11):
        Bayesian nonparametrics and practical applications.
\end{enumerate}

\section*{Pedagogical philosophy}

Each chapter alternates between:
\begin{itemize}
  \item rigorous \textbf{theoretical developments} (definitions,
        theorems, proofs),
  \item detailed \textbf{worked examples} with explicit calculations,
  \item \textbf{numerical implementations} in Python (PyMC, NumPy,
        SciPy),
  \item \textbf{exercises} of increasing difficulty.
\end{itemize}

\section*{Bayesian \textit{versus} frequentist approaches}

Before delving into the detailed study, let us contrast the two
paradigms:

\begin{center}
\renewcommand{\arraystretch}{1.3}
\begin{tabular}{p{4cm} p{5cm} p{5cm}}
\toprule
\textbf{Aspect} & \textbf{Frequentist} & \textbf{Bayesian} \\
\midrule
Parameter $\theta$ & Fixed, unknown & Random variable \\
Probability & Limiting frequency & Degree of belief \\
Uncertainty & Confidence intervals & Credible intervals \\
Prior information & Not formally used & Integrated via prior \\
Main result & Estimator $\hat{\theta}(x)$ & Posterior $\pi(\theta \mid x)$ \\
Coherence & May violate likelihood principle & Coherence guaranteed \\
Small samples & Possible difficulties & Naturally suited \\
Computation & Often analytic & Often numerical (MCMC) \\
\bottomrule
\end{tabular}
\end{center}

\section*{The founding theorem}

The entire Bayesian paradigm rests on a single formula:

\begin{keyformulas}[title=\textbf{Bayes' Theorem}]
Let $\theta \in \Theta$ be a parameter with prior distribution
$\pi(\theta)$ and $x = (x_1, \ldots, x_n)$ a sample with likelihood
$L(\theta \mid x)$. The \textbf{posterior distribution} is:
\[
  \pi(\theta \mid x)
  = \frac{L(\theta \mid x)\,\pi(\theta)}
         {\int_{\Theta} L(\theta \mid x)\,\pi(\theta)\,d\theta}
  \propto L(\theta \mid x)\,\pi(\theta).
\]
\end{keyformulas}

\section*{Notation and conventions}

Throughout this course, we adopt the following conventions:

\begin{itemize}
  \item $X$ denotes a random variable, $x$ an observed value.
  \item $\theta$ is the parameter of interest, $\Theta$ the parameter
        space.
  \item $\pi(\theta)$ is the prior distribution.
  \item $\pi(\theta \mid x)$ is the posterior distribution.
  \item $L(\theta \mid x) = f(x \mid \theta)$ is the likelihood.
  \item $m(x) = \int L(\theta \mid x)\,\pi(\theta)\,d\theta$ is the
        marginal likelihood (evidence).
  \item $\E[\cdot]$, $\Var[\cdot]$, $\Cov[\cdot]$ denote
        expectation, variance, and covariance.
  \item $\PP(\cdot)$ denotes probability.
  \item $\KL(p \| q)$ denotes the Kullback--Leibler divergence.
  \item $\norm{\cdot}$ denotes the Euclidean norm.
\end{itemize}

\section*{Commonly used distributions}

\begin{center}
\renewcommand{\arraystretch}{1.3}
\begin{tabular}{lll}
\toprule
\textbf{Distribution} & \textbf{Notation} & \textbf{Parameters} \\
\midrule
Bernoulli & $\text{Ber}(p)$ & $p \in [0,1]$ \\
Binomial & $\text{Bin}(n, p)$ & $n \in \N$, $p \in [0,1]$ \\
Poisson & $\text{Poi}(\lambda)$ & $\lambda > 0$ \\
Normal & $\mathcal{N}(\mu, \sigma^2)$ & $\mu \in \R$, $\sigma^2 > 0$ \\
Gamma & $\text{Ga}(\alpha, \beta)$ & $\alpha, \beta > 0$ \\
Beta & $\text{Be}(a, b)$ & $a, b > 0$ \\
Dirichlet & $\text{Dir}(\boldsymbol{\alpha})$ & $\alpha_k > 0$ \\
Student-$t$ & $t_\nu(\mu, \sigma^2)$ & $\nu > 0$ \\
Wishart & $\mathcal{W}_p(\nu, \boldsymbol{\Sigma})$ & $\nu > p-1$ \\
Inverse-Gamma & $\text{IG}(\alpha, \beta)$ & $\alpha, \beta > 0$ \\
\bottomrule
\end{tabular}
\end{center}

\section*{Historical milestones}

\begin{itemize}
  \item \textbf{1763}: Posthumous publication of Thomas Bayes' theorem
        by Richard Price.
  \item \textbf{1812}: Pierre-Simon de Laplace formalizes and generalizes
        the method in \textit{Th\'eorie analytique des probabilit\'es}.
  \item \textbf{1937}: Bruno de Finetti develops the notion of subjective
        probability and the exchangeability theorem.
  \item \textbf{1954}: Leonard Jimmie Savage publishes \textit{The
        Foundations of Statistics}, a pillar of Bayesian decision theory.
  \item \textbf{1970s}: Dennis Lindley and others establish the modern
        foundations.
  \item \textbf{1990s}: The MCMC revolution --- Metropolis--Hastings
        becomes the workhorse of Bayesian computation.
  \item \textbf{2000s--2020s}: Variational inference, Gaussian processes,
        Bayesian deep learning.
\end{itemize}

\section*{Software and libraries}

The practical implementations in this course use:

\begin{itemize}
  \item \textbf{Python 3.10+} as the primary language.
  \item \textbf{PyMC} (v5+) for Bayesian modeling and MCMC.
  \item \textbf{ArviZ} for chain diagnostics and visualization.
  \item \textbf{NumPy / SciPy} for numerical computation.
  \item \textbf{Matplotlib / Seaborn} for plotting.
\end{itemize}

\begin{codeblock}[title=\textbf{Installing dependencies}]
\begin{minted}{python}
# Recommended installation
# pip install pymc arviz numpy scipy matplotlib seaborn

import pymc as pm
import arviz as az
import numpy as np
import matplotlib.pyplot as plt

print(f"PyMC version: {pm.__version__}")
print(f"ArviZ version: {az.__version__}")
\end{minted}
\end{codeblock}

\section*{How to use this course}

\begin{warning}[title=\textbf{Important prerequisites}]
This course assumes that the student is proficient in:
\begin{enumerate}
  \item Probability theory (random variables, classical distributions,
        modes of convergence).
  \item Mathematical statistics (estimation, testing, likelihood).
  \item Linear algebra (matrices, decompositions, quadratic forms).
  \item Python programming (basic syntax, NumPy).
\end{enumerate}
\end{warning}

\begin{intuition}[title=\textbf{Guiding thread}]
The guiding idea of this course is simple: \textit{all inference is a
problem of updating a probability distribution}. We start from a prior
belief $\pi(\theta)$, we observe data $x$, and we obtain a posterior
belief $\pi(\theta \mid x)$. This unified vision underlies all the
methods presented herein.
\end{intuition}

\section*{Main references}

\begin{itemize}
  \item A.~Gelman, J.B.~Carlin, H.S.~Stern, D.B.~Dunson, A.~Vehtari,
        D.B.~Rubin, \textit{Bayesian Data Analysis} (3rd ed.), 2013.
  \item C.P.~Robert, \textit{The Bayesian Choice} (2nd ed.), 2007.
  \item J.K.~Kruschke, \textit{Doing Bayesian Data Analysis} (2nd ed.),
        2015.
  \item D.~Barber, \textit{Bayesian Reasoning and Machine Learning}, 2012.
  \item K.P.~Murphy, \textit{Machine Learning: A Probabilistic Perspective},
        2012.
  \item P.~Hoff, \textit{A First Course in Bayesian Statistical Methods},
        2009.
\end{itemize}

\bigskip
\begin{flushright}
\textit{The author} \\
March 2026
\end{flushright}
